\section{Introduction}
\label{sec:introduction}

Phase interpolators (PIs) are a critical building block in high-speed serial
links and mm-wave clock generation circuits, providing fine-grained phase
selection between a set of reference phases. At millimeter-wave frequencies
(\SIrange{30}{300}{\giga\hertz}), the design of PIs is challenged by the
reduced transconductance efficiency of transistors, increased sensitivity to
parasitic capacitances, and tight timing margins.

A common technique to ease these constraints is frequency doubling: the
phase interpolation is performed at a sub-harmonic (typically \(f_0/2\)) and
the output is then frequency-doubled to the target mm-wave frequency. This
relaxes the bandwidth requirement on the interpolating core and allows the
use of lower-speed, more linear combining topologies.

The remainder of this paper is organized as follows.
Section~\ref{sec:architecture} describes the top-level architecture.
Section~\ref{sec:circuit} details the transistor-level implementation.
Section~\ref{sec:analysis} provides a theoretical analysis of phase
resolution and nonlinearity. Section~\ref{sec:simulation} presents simulation
results, and Section~\ref{sec:conclusion} concludes the paper.
